\chapter{Background}
\label{chap:Multicat}
\section{Multicategories}
\label{sec:multicat}
{\bf EXPLAIN THE SOURCE OF THIS MATERIAL. ``Multicategories were first defined in ...'' ``Our exposition follows ...''}
yes
Multicategories are a generalization of categories. In 
categories, arrows have a single source/domain to a single target/codoamin, in a multicategories the arrows have a finite sequence of sources/inputs to a single target/output. 
\begin{defn}
    A \emph{multicategory} $\M$ consist of a set $S$ of objects and for each $n\geq 0$ and each sequence $s_1,\ldots, s_n,t$ of objects a set
    \[\M(s_1,\ldots,s_n;t) \] 
    of multiarrows/operations, thought of as taking $n$ inputs of objects $s_1,\ldots,s_n$ respectively to the object output $t$. Moreover, there are three kinds of structure maps:
    \begin{itemize}
        \item For each object $s$, an identity/unit $1_s\in \M(s;s)$
        \item For any sequence of objects $s_1,\ldots,s_n,t$ and any $n$-tuple of sequences $d^i_1,\ldots,d^i_{k_i}$ of objects for $i=1,\ldots,n$, a composition function of multiarrows/operations 
        \[
        \gamma:\M(s_1,\ldots,s_n;t)\times\prod_{i=1}^{n}\M(d^i_1,\ldots,d^i_{k_i};s_i)\rightarrow\M(d^1_1,\ldots,d^1_{k_1},d^2_1,\ldots,d^n_{k_n};t),
        \]
        usually written $\gamma(p,q_1,\ldots,q_n)=p\circ(q_1,\ldots,q_n)$ or even just $p(q_1,\ldots,q_n)$
    \end{itemize}
    Furthermore these structure maps have to satisfy a number of axioms, which express that composition is unital and associative \citep{HM}
    % {\bf ADD SPECIFIC CITATION WHERE THE FULL DEFINITION APPEARS.}
\end{defn}

\begin{rmk}
Let $\M$ be a multicategory, $f\in\M(s_1,\ldots,s_n;t)$ and $g\in\M(d_1,\ldots,d_l;s_i)$ for some $1 \leq i \leq n$. It will be convenient to have the following notation:
$$
\iof{f}{g} := f(1_{s_1},\ldots,g,\ldots,1_{s_n})
$$
Thus, $\iof{f}{g}$ is the composition of $f$ with $g$ in the $i$th variable.
\end{rmk}

\begin{ex} 
    You can think of the multiarrows in $\M(s_1,\ldots, s_n;t)$ as operations, taking inputs of types $s_1,\ldots,s_n$ respectively to an output of type $t$. 
     {\bf FIX A SINGLE SET $X$? and then $n$-ary multiarrows are operations on $X$.}
    For example, for a fixed family of sets $ {X_c}_{c\in C}$, define a multicategory $\M$ whose objects is the set of colors $C$,  and whose multiarrows are defined by  $\M(c_1,\ldots, c_n;c)$ to be the set of functions 
     \[
     p:X_{c_1}\times \cdots \times X_{c_n}\rightarrow X_c
    \]
    Composition in $\M$ is defined as composition of functions.
\end{ex}


\begin{ex}
    If $\C$ is a symmetric monoidal category and $C$ is a set of objects of $\C$, then one obtains a multicategory with this set of objects by defining $\M(c_1,\ldots,c_m;t)$ to be the set of morphisms $c_1\otimes \dots \otimes c_n \rightarrow t$ in $\C$. We will sometimes write $\C^{\otimes}$ for the multicategory obtained in this way, when $C$ is the set of all objects of $\C$ (granted it is small).
\end{ex}
\begin{ex}
    A (small) category $\C$ can be considered as a multicategory $\C_m$ which only have unary arrows. I.e. the set $\C_m(c_1,\ldots,c_n;c)$ is always empty unless $n=1$.
\end{ex}

We can observe that multicategories from a category $\MCat$.


\begin{defn}
For two multicategories $\M$ and $\mathcal{N}$, a \emph{multifunctor} $\phi:\M\rightarrow\mathcal{N}$ consist of:
\begin{itemize}
    \item A function $f:O_{\M}\rightarrow O_{\mathcal{N}}$ between the corresponding sets of objects
    \item For each sequence $s_1,\ldots,s_n,t$ of objects of $\M$, a map 
    \[
        \phi_{\lst{s}{n},t}:\M(\lst{s}{n};t)\rightarrow \mathcal{N}(\flst{s}{n}{f};f(t)).
    \]
\end{itemize}
These maps being compatible with compositions and units in the obvious way.
\end{defn}
\subsection{Trees}
\label{sec:trees}

\begin{defn}
    A \emph{tree} $T$ consist of a finite set $V$ of \emph{vertices}, a nonempty finite set $E$ of \emph{edges}, a distinguished element $r\in E$ called the \emph{root}, together with the following data:
    \begin{enumerate}[label=(\arabic*)]
        \item a function $I:E-\{r\}\rightarrow V$, which we think of as assigning to an edge $e$ the vertex $I(e)$ of which is an input.
        \item A function $O:V\rightarrow E$, assigning to each vertex $v$ its output edge $O(v)$
    \end{enumerate}
\end{defn}

For each edge $e$ other than the root $r$, we obtain a sequence of edges starting at $e$ by repeatedly applying $O\circ I$. We demand that this sequence ends in the root $r$ after finitely many steps, for an arbitrary starting edge.
The edges in the complement of the image of $O$ are called \emph{leaves} of the tree $T$. The vertices in the complement of the image of $I$ are called \emph{stumps}, or \emph{nullary vertices.} An \emph{outer edge} is an edge tha is either the root or a leaf. An \emph{inner edge} is any other edge of $T$. Such an edge is both an output edge and an input edge for some other vertex.\\

When writing $T$ for a tree, we will usually write $E(T)$ and $V(T)$ for its sets of edges and vertices, repectively.

\begin{ex}
    The smallest possible tree consiste of a single edge and no vertices; this tree will be denoted by $\mu$
\end{ex}
\begin{ex}
The next smallest tree consist of a single edge and single vertex, pictured as:
% https://q.uiver.app/#q=WzAsMyxbMCwxLCJcXGJ1bGxldCJdLFswLDJdLFswLDBdLFsyLDEsIiIsMCx7InNob3J0ZW4iOnsic291cmNlIjo1MH0sInN0eWxlIjp7ImhlYWQiOnsibmFtZSI6Im5vbmUifX19XV0=
\[\begin{tikzcd}
	{} \\
	\bullet \\
	{}
	\arrow[between={0.5}{1}, no head, from=1-1, to=3-1]
\end{tikzcd}\]
\end{ex}
\newpage
\begin{ex}
    A bigger and more typical tree $T$ with $E(T)=\{x,y,z,w,v,r\}$ and $V(T)=\{f,g,h,l\}$:\\
    \begin{center}
\begin{tikzpicture}[
  vertex/.style={circle, fill=black, inner sep=1.5pt},
  edge/.style={draw},
  lab/.style={font=\small}
]

% vertices (explicit coordinates)
\node[vertex, label=above:$f$] (f) at (-.2,0) {};
\node[vertex, label=right:$g$] (g) at (-.7,-1.4) {};
\node[vertex, label=left:$l$] (l) at (.6,-3.0) {};
\node[vertex, label=above:$h$] (h) at (.6,-1.4) {};

% dangling endpoints (no vertices)
\coordinate (xend) at (-1.3,-.2);
\coordinate (wend) at (1.9,-1.5);
\coordinate (rend) at (0.6,-4.6);

% edges
\draw[edge] (xend) -- node[lab, left, pos = .45] {$x$} (g);
\draw[edge] (f) -- node[lab, right] {$y$} (g);

\draw[edge] (g) -- node[lab, left] {$v$} (l);
\draw[edge] (h) -- node[lab, right] {$z$} (l);

\draw[edge] (l) -- node[lab, right ] {$w$} (wend);
\draw[edge] (l) -- node[lab, left] {$r$} (rend);

\end{tikzpicture}
\end{center}
\end{ex}



% \begin{thm}
% Something Ramdon

% \end{thm}

% \begin{defn}
% definition now. Here is $\Set$
% \end{defn}

% \begin{rmk}
%     here goes it
% \end{rmk}